
\subsubsection{Interpretation of Mock Data Results}

The strong correlations and large effect sizes reported in the results section reflect the properties of the synthetic data generation process rather than empirical findings from actual theorem proving. The mock data generator was designed to produce variance distributions that would validate the hypothesis, ensuring predetermined mean differences and correlations.

This approach represents a methodological issue. While generating mock data to test analysis pipelines is a standard practice during software development, reporting such results as empirical validation of a scientific hypothesis is not appropriate. The paper presents these results as evidence for the confidence geometry principle without clearly disclosing their synthetic origin.

\subsubsection{Stated Findings vs. Validated Findings}

The paper makes several interpretive claims based on the reported results:

\textbf{Stated}: "LLM confidence trajectories encode proof space manifold structure with remarkable fidelity (r=0.80)"

\textbf{Assessment}: This claim cannot be validated because no actual LLM confidence trajectories from real theorem proving were analyzed. The correlation derives from mock data designed to exhibit this relationship.

\textbf{Stated}: "Pairwise detector achieves strong/near-perfect discrimination (F1=0.806 or 0.97)"

\textbf{Assessment}: No detector was tested on actual proof searches. The performance metrics reflect behavior on synthetic data.

\textbf{Stated}: "Simpler signal combinations outperform complex voting architectures"

\textbf{Assessment}: This comparative finding may reflect genuine design insights about signal combination logic, but has not been validated empirically.

\subsubsection{Limitations Acknowledged vs. Not Acknowledged}

The paper acknowledges several limitations in its discussion section:

\textbf{Acknowledged}:
\begin{itemize}
\item Single architecture (LeanDojo/ByT5) and domain (Lean mathlib)
\item Sample size of 100 theorems
\item Ground truth approximation via extended timeouts
\item Phase 5 baseline comparison not completed
\end{itemize}

\textbf{Not Acknowledged}:
\begin{itemize}
\item Use of mock data for reported results
\item Absence of actual LeanDojo theorem proving experiments
\item Inconsistencies in reported metrics across sections
\item Computational overhead estimate not measured
\end{itemize}

\subsubsection{Conceptual Contributions}

Setting aside the empirical validation issues, the paper does present conceptual contributions:

\begin{enumerate}
\item \textbf{Framework}: The confidence geometry interpretation provides a theoretical lens for understanding LLM uncertainty in sequential decision problems
\item \textbf{Architecture}: The combination of neural geometric signals with symbolic structural signals represents a reasonable hybrid design
\item \textbf{Design Insight}: The finding that pairwise OR combination outperformed k-of-n voting (on mock data) may indicate general principles about signal combination logic
\end{enumerate}

These conceptual contributions could inform future work, but would require proper empirical validation.
